\subsection*{Proposition~\ref{prop:nonid}: the two-point algebra in full}

Let a scoring rule take one of two values, with $a$ of the $P$ positives and
$b$ of the $N$ negatives in the high block; write $\alpha = a/P$,
$\beta = b/N$, $\rho = N/P$ and $\pi = P/(P+N) = 1/(1+\rho)$. Two exact
identities follow directly from the definitions in
Section~\ref{sec:props}:
\[
\mathrm{AUC} = \tfrac12 + \tfrac{\alpha - \beta}{2},
\qquad
\mathrm{AP} = \alpha\,\frac{\alpha}{\alpha + \beta\rho} + (1-\alpha)\pi .
\]
The first is the Mann--Whitney statistic with ties at $1/2$; the second is the
step-wise sum over the two operating points the rule admits. Hence
\[
\mathrm{AP} - \pi \;=\; \alpha\left(\frac{\alpha}{\alpha + \beta\rho} -
\pi\right).
\]

Fix $d = \alpha - \beta > 0$. Every rule on the line $\beta = \alpha - d$ has
the same AUC increment $d/2$ over the fully tied baseline, so any two of them
give $R_{\mathrm{AUC}} = 0$ exactly. Along that line the AP increment runs
from $d(1-\pi)$ at $\alpha = d$, $\beta = 0$ down to
$\dfrac{d\rho}{\bigl(1 + (1-d)\rho\bigr)(1+\rho)}$ at $\alpha = 1$, a ratio of
about $1 + \rho$. The construction therefore needs an imbalanced problem to
drive $R_{\mathrm{AP}}$ close to one, and the file uses $\rho = 500$. This is
a property of the construction and not a hidden assumption: with $\rho = 9$
the same argument gives a gap of about $0.9$, which is still enough to refute
the existence of $\varphi$.

\subsection*{Proposition~\ref{prop:rank}: the cancelling configuration}

Remark~\ref{rem:cancel} claims a configuration on which a non-rank-based
metric's increment moves under $\psi$ while the ratio does not. The
construction is immediate once stated: take $\sigma_{B_1} = \sigma_{B_0}$ and
$\sigma_{B_1 \cup f} = \sigma_{B_0 \cup f}$. Then $V(B_1) = V(B_0)$ for every
metric and every $\psi$, so $R = 0$ identically, while $V$ itself moves under
$\psi$ whenever $m$ is not rank-based. The configuration is degenerate --- the
two baselines induce the same scores --- and that is the point: the
existential quantifier in part~(ii) cannot be strengthened to a universal one
without excluding cases like it, and the earlier version of this work did not
exclude them.

\subsection*{Result~\ref{res:nonredundancy}: what the search does}

For each ordered pair of axes $(i,j)$ the search looks for two estates whose
increments agree, to a declared tolerance, at \emph{every} level of axis $i$
and disagree at some level of axis $j$. A reader who knew only the axis-$i$
readings could not tell those two estates apart, yet their axis-$j$ readings
differ; so the axis-$i$ readings do not determine the axis-$j$ readings.

The family is \nResultEstates\ synthetic estates, each with its own seed,
prevalence, signal strength and noise; every cell is a pure function of the
estate and the level assignment. An earlier implementation drew fresh noise
inside each cell evaluation, so two evaluations of the same cell disagreed and
the search found almost nothing; that is a defect of the search, not a
property of the axes, and it is recorded here because the count it produced
appeared in a previous draft.
